% !TEX root = ../UPDATED_main.tex
\section*{Discussion}

SatuTe asks whether an alignment retains detectable dependence between the subtrees separated by a branch. This question differs from bootstrap and likelihood-ratio support, which quantify the reproducibility of, or preference for, a split \cite{felsenstein1985confidence,hoang2017,lemoine2018}. The original JC benchmark makes the distinction explicit: bootstrap support can remain complete after SatuTe reaches its null rejection rate. A non-significant SatuTe result therefore indicates insufficient branch-specific signal in the tested alignment and model; it does not reject common ancestry.

The original SatuTe formulation remains the primary model in this study. It supplies the null hypothesis, branch-specific interpretation, one-sided decision rule, and baseline simulation design. The eigenvalue-weighted statistic is a model-aware extension that preserves each of these elements. It reduces exactly to original SatuTe under JC but allows faster non-stationary directions to contribute while they persist under GTR or amino-acid models.

The benefit of weighting was localized rather than universal. Gains occurred mainly in short alignments near the long-branch detection boundary and were negligible when both formulations had high power or had reached the null floor. Branch-neighborhood audits showed that weighting did not indiscriminately increase scores across a tree. The size of the effect also depended strongly on the substitution matrix: JTT produced broader shifts than LG or Q.pfam, showing that alphabet size alone does not determine the available spectral information.

Model adequacy was at least as consequential as the choice of statistic. JC and K2P analyses of GTR data were too liberal, whereas F81 was closer to the correctly specified result. The JC zero-difference control establishes equality between the original and weighted formulas under JC; it does not establish the adequacy of JC for data generated under a different process.

Invariant sites create a separate structural problem. A zero-rate category remains ancestrally dependent across an infinite focal branch. Placing that category in a product-form null therefore introduces a positive mean that studentization amplifies as alignment length increases. Including the invariant category's joint probability only in the null denominator restored exact centering and standard-normal calibration. After independent empirical thresholding, gains again followed the detection boundary. The exploratory \(+I+R4\) analysis serves as a stress test, whereas fixed GTR+F+I+G4 is the publication benchmark.

The empirical examples demonstrate the value of a local diagnostic. Ribosomal proteins retained signal across nearly all windows, whereas 16S rRNA contained broad saturated regions and lost signal for the Eukaryota branch in its fastest category. In HCV, whole-gene tests were uniformly informative, yet short-window losses concentrated at synonymous-enriched third positions. An alignment-wide pass can therefore conceal within-locus saturation, whereas SatuTe profiles can guide partition-specific model checks or selective masking without discarding an entire coding region. Weighting reduced the number of threshold crossings without changing either biological contrast.

These examples also define the limits of interpretation. SatuTe localizes signal but does not select the correct topology, establish model adequacy, or estimate viral population diversity. The HCV analysis is further limited by its ten-genome reference panel, overlapping short windows, and model refitting at that scale. Non-stationary or non-reversible extensions would make the root identifiable and require a root-aware null.

Within the present reversible framework, we recommend reporting original SatuTe as the primary branch-specific diagnosis, accompanied where appropriate by weighted sensitivity results, paired differences, multiplicity correction, conventional branch support, and model-adequacy checks. Integrated into IQ-TREE, this reporting structure makes the loss of historical signal testable at the level where phylogenetic conclusions are drawn: individual branches and alignment regions.
